\section{Implementation details of the frozen prospective study}
\label{sec:implementation-details}

\PaperTable{tables/backend_spec.tex}

% Implementation bindings: I01--I03. Code/protocol inspection; no execution.
\subsection{Prompt actions, candidate sets and STOP}
The frozen vocabulary uses the VOC canonical class strings, also on COCO:
\emph{aeroplane, bicycle, bird, boat, bottle, bus, car, cat, chair, cow,
dining table, dog, horse, motorbike, person, potted plant, sheep, sofa,
train, tv/monitor}. A0 applies the canonical text. A1 lowercases and
normalizes whitespace, then uses the fixed taxonomy mapping; the normalized
strings in this vocabulary coincide with the canonical strings. A2 uses
\texttt{a photo of a \{normalized class\}}. These remain separately cached
registered calls; equal strings do not authorize dropping an action.

A3 applies the canonical prompt after exact horizontal pixel reversal and
reverses predicted masks back to the original coordinates. A4 crops around
one A0 source candidate: expand its half-open bounding box by factor 1.25
about its center, floor lower bounds, ceil upper bounds, and clip to the
image. Apply the canonical prompt in that crop, restore predictions inside
the crop window, and set outside pixels false. There is no project-side
resize in addition to the frozen backend preprocessing. A5 supplies the
source box as a positive box prompt, converted to normalized XYWH. A6
chooses the pixel of maximum four-connected erosion depth inside the source
mask, breaking ties by row-major order. Its positive point is converted
using the pixel center; a fresh session replays the canonical prompt and
refines the source's upstream object ID. A6 requires that ID in addition to
the source mask.

The source is the A0 candidate with greatest upstream \texttt{out\_probs},
with ascending candidate ID breaking ties. This planning choice differs
from the outcome evaluated for each action: the output is the union of all
its retained candidates. In particular, STOP uses the A0 union, not the
source candidate, and incurs no additional call or cost. The frozen adapter
consumes the upstream binary masks; it does not tune a new segmentation
threshold for this study. It preserves nonempty candidates and their raw
indices, while filtering empty candidates before candidate-record creation.

\subsection{Outcome validity is distinct from structural feasibility}
A completed, structurally valid response with no retained candidates has a
defined all-false union. It is a legal no-result, not a zero-loss target:
its mismatch still depends on the reference label. When A0 has this outcome,
STOP is defined, A3 remains independently feasible, and A4--A6 are
structurally infeasible because no A0 source exists. An unavailable upstream
object ID independently makes A6 infeasible. Infeasible actions are masked
from decisions and have no observed counterfactual target.

A technical failure supplies neither an empty mask nor a class-absence
statement. The frozen external missingness disposition preserves the panel
state but leaves the affected target undefined. For the trajectory-level
quarantine, prediction is also unavailable. On the three A4-failure states,
the original plan still marks A4 structurally feasible. After the failures
are recorded, the pre-GT prediction-lock pipeline copies the plan mask and
sets these A4 entries false for offline selection. The stored ensemble field
named \texttt{feasibility} is this adjusted availability mask; it does not
redefine the original action-plan feasibility or change the continuous
predictions. Being fixed before GT access makes the adjustment label-blind,
but the failure information is acquired after action execution and is not
available to the initial repair decision.

The evaluator uses original plan feasibility and target definedness. It
retains unaffected fidelity cells and defined STOP-relative pairs, while
excluding each affected state from complete-oracle regret, realized-gain
and gain-capture summaries, and both sparse-budget populations. The
defined-prediction non-STOP diagnostic instead includes these states in its
counting population; only non-STOP choices contribute to the count retained
by the frozen E4-C predicate. These offline
choices are not evidence for a deployable failure-aware policy. No failure
is imputed as STOP, zero gain, or a valid empty response, and no failed
repair is relabeled structurally infeasible.

% Implementation bindings: I04--I06.
\subsection{Candidate-level cells and matched region controls}
Atom construction receives every retained candidate from A0, A1 and A2.
Candidates are ordered by mask digest, candidate ID and action ID; distinct
candidate records retain their identity even when masks coincide. The
membership signature has one coordinate per candidate. Four-connected
components of identical signatures form the cells, including the zero
signature outside every candidate. There is one full-raster cell if the
candidate list is empty. Stable component ordering is followed by explicit
label-index remapping. These details preserve candidate-level incidence,
which would be lost by first replacing each action's candidate set with its
union. At F1, observed A3 candidates enter the same construction; A4--A6
remain excluded.

Rect starts with the full raster and recursively splits rectangles until
their count equals the Atom cell count. At each step it selects a largest
positive-area rectangle, preferring topmost then leftmost position on ties.
It prefers the longer axis (horizontal extent on a tie), searches cuts with
two positive areas, and chooses the smallest area imbalance with the earlier
cut breaking ties. A second axis is considered if necessary. Final region
ordering is deterministic. This matches region count, not each Atom area's
value. The frozen identifier ``area-matched rectangles'' is retained in the
artifact registry as a historical name for this algorithm.

Union forms the union of all observed candidates and separately decomposes
that union and its complement into four-connected components. If the union
is all false or all true, there is one full-raster component. The stored
bounding boxes summarize components; they are not overlapping rectangular
prediction supports. Atom, Rect and Union all supply disjoint exhaustive
partitions, so each pixel receives one region assignment. Each model's
prediction weights are its own full-raster region-area fractions. Target
validity can remove supervision from a region without removing its feature
geometry; the target-only reconstruction weights remain separate.

\subsection{Numeric features and their information content}
The prospective region vector has 27 coordinates: area fraction; normalized
bounding box and centroid; logarithms of component index and adjacency
degree; normalized shared-boundary length; canonical-union membership,
disagreement and consensus flags; A0/A1/A2 availability flags; logarithmic
signature length and mean membership; logarithmic incidence count; and an
eight-bin signed candidate-incidence hash. The canonical-union bit indicates
membership in at least one A0 candidate. The signature is summarized into
fixed-dimensional features rather than supplied as a variable-length binary
network input.

The 41 state coordinates contain logarithmic region count, five area-fraction
summaries (minimum, median, mean, maximum and standard deviation), four-way
status encodings for each base action, disagreement/conflict fractions,
logarithmic already-incurred base cost, and the 20-way class indicator.
The 16 action coordinates contain feasibility, source presence, five-way
action identity, precondition availability, and an eight-bin signed hash of
the source ID and plan lineage. The old realized-runtime coordinate is
absent. Hash encodings use fixed namespaces and deterministic SHA-256 bins
and signs, normalized by the square root of item count; they add no learned
image embedding. G-small, G-match and Gain concatenate the state vector
with one action vector; Defer and SPO+ concatenate the state with all five
action vectors in physical order.

The region controls preserve dimensions but not every descriptor value.
Rect has empty signature/incidence fields and zero canonical-membership,
disagreement and consensus features. Union uses an inside/outside signature
and the synthetic incidence token \texttt{OBSERVABLE\_UNION}; its canonical,
disagreement and consensus flags are also zero. Each has its own region
statistics, and Union records union coverage in its state descriptor. The
F1 encoding includes acquired A3 status through a fixed observation-context
token in the incidence hash/count and adds acquired A3 cost to base cost.
Thus representation comparisons share model dimensions and supervision,
while their documented feature content follows the chosen geometry and
filtration.

% Implementation bindings: I07--I09.
\subsection{Network graphs and output semantics}
G-small is a $57\!\to\!128\!\to\!1$ multilayer perceptron, with LayerNorm
and SiLU after the hidden affine layer, dropout probability 0.10, and a
sigmoid scalar output. G-match adds a second hidden layer and uses width
292: $57\!\to\!292\!\to\!292\!\to\!1$, with LayerNorm on the first
hidden layer, SiLU on both, dropout after the first, and a final sigmoid.
The width is selected from integers 1--4096 by parameter-count distance to
Atom, with the smaller width breaking a tie; no targets enter this search.

In the shared local graph, a $27\!\to\!128\!\to\!128$ region encoder
uses LayerNorm after its first affine layer and SiLU after both. Concatenate
the unweighted region-embedding mean and maximum with the 41 state
coordinates, then apply a $297\!\to\!128$ LayerNorm--SiLU context layer.
A learned five-entry, 32-dimensional action-ID embedding concatenates with
the 16 action coordinates and passes through a $48\!\to\!64$ SiLU layer.
Each region/action pair concatenates its 128-dimensional region embedding,
128-dimensional state context and 64-dimensional action embedding. A
$320\!\to\!128\!\to\!1$ head uses LayerNorm, SiLU, dropout 0.10 and a
sigmoid. Its regional outputs aggregate with observable area fractions.
Atom, Rect, Union, Regret, Relative and F0/F1 share this graph; the latter
variants differ in their registered objective, representation or observation.

Gain uses $57\!\to\!128\!\to\!128\!\to\!1$ with the same first-layer
normalization/dropout and hidden SiLU pattern, but an unbounded linear
output. Its scientific STOP gain is zero after seed averaging. Defer and
SPO+ share $121\!\to\!128\!\to\!5$, with LayerNorm, SiLU and dropout
0.10, and an unbounded output vector. Defer chooses the feasible maximum
logit; SPO+ chooses the feasible minimum value. Their dispatch scores are,
respectively, best non-STOP minus STOP logit, and STOP minus best non-STOP
predicted value. No residual MAE is assigned to these decision outputs.
The canonical parameters and alias identities are given in
Table~\ref{tab:implementation-panel}.

% Implementation bindings: I10--I12.
\subsection{Cross-fitting, scaling and optimization}
The fixed five-fold manifest groups every image's class states together.
Its original construction uses capacity-constrained multilabel balancing
with deterministic seeded ordering; it is reused, not regenerated per
model. Within an outer fold and seed, candidate training groups are ordered
by a hash of scale, seed, outer-fold identity and image-group ID. The first
$\lfloor n_{\rm outer\text{-}train}/5\rfloor$ groups form inner-stop;
the remainder supply optimization losses. This is repeated for seeds 13,
37 and 71 with the same registered outer folds.

State, action and region scalers are separate. Means and population standard
deviations are computed in float64 from outer-training groups, including
inner-stop but excluding the outer-held-out groups. Standard deviations
below $10^{-8}$ are replaced by one; mean/scale and transformed vectors
are stored in float32. Action scaling follows removal of the runtime
coordinate and includes the five physical plan rows. Region scaling uses
all regions belonging to the same training groups. The F0/F1 action scaler
instead fits and transforms only STOP/A4/A5/A6; the physical A3 row is
initialized to exact zeros and contributes to neither scaling nor loss.

V6's first batch, decision-objective batch and F0/F1 use AdamW with learning
rate $5\times10^{-4}$ and weight decay $10^{-4}$. Each epoch is one full
optimization-partition gradient step; there is no minibatch-size parameter
to infer from an artifact name. The full feature surface can be forwarded
while the loss mask selects only training states; region/context reductions
are within state. Activation checkpointing/chunking bounds large-region
memory without changing the reduction. Early stopping monitors the same
registered model-specific objective on inner-stop groups, retains a
checkpoint only for a decrease greater than $10^{-8}$, and uses patience
15 with a maximum of 120 epochs. First-batch and F0/F1 objectives are
state-only Huber; the decision-family objectives are given in
Appendix~\ref{sec:accounting-appendix}. NumPy and PyTorch seeds are fixed,
deterministic PyTorch operations are requested, cuDNN benchmarking and TF32
are disabled, and the prescribed CUDA workspace setting is retained.

\subsection{Full-fit transfer and ensemble order}
The external checkpoints are fresh full fits on all authorized S1364 FIT
groups. Each family's epoch count is the eighth sorted value among its 15
registered V6 best epochs, shared by all three seeds. The frozen resulting
counts are 120 for G-small, G-match, Atom, Rect and Union, and 108 for
Relative. In fact, all 15 saved best epochs for each of those first five
families are 120, the development maximum. Relative's sorted sequence is
29, 59, 69, 84, 96, 98, 101, 108, 111, 120, 120, 120, 120, 120, 120.
These fixed-budget comparisons establish neither convergence nor
undertraining or ranking stability under further optimization.
There is no external-label early stopping, warm start, scheduler,
mixed-precision or distributed-data-parallel training in that refit. The
final fixed-epoch checkpoint is retained. Full-fit scalers use all FIT rows;
COCO features only apply those stored scalers and never fit new statistics.

Development stitches outer-held-out predictions per seed before arithmetic
averaging. External transfer averages the three continuous seed surfaces in
float64 before action selection. Action ties use STOP/A3/A4/A5/A6 order;
sparse state ties use ascending semantic state ID. F0/F1 retain only their
registered remaining actions. The ensemble averages predictions, not
parameters or votes.
